% pset mechanism, generated by pset PSET_VERSION.
%
% Do not edit: `pset --init` overwrites this file to upgrade it.  Styling
% belongs in your own preamble, through the hooks below.
%
% \sol and \rub print only in the build that reveals them, \stud only in
% the build that reveals neither.  pset defines \showsol and \showrub to
% pick the build, deliberately not \sol and \rub: a document cannot both
% define a macro and test whether it is defined, and that collision
% prints the answers on the student copy.

% Rendering hooks, plain by default so this file needs no packages and
% cannot collide with your package set.  \renewcommand any of them after
% inputting this file to style the result.
\providecommand{\solstyle}[1]{#1}
\providecommand{\rubstyle}[1]{#1}
\providecommand{\probstyle}[2]{\bigskip \par \noindent { \textbf{Problem #1 #2} \\ }}

\newcommand{\sol}[1]{}
\newcommand{\rub}[1]{}
\newcommand{\stud}[1]{#1}

% answer key: show \sol, drop the student-only \stud
\ifdefined\showsol
	\renewcommand{\sol}[1]{\solstyle{#1}}
	\renewcommand{\stud}[1]{}
\fi

% rubric copy: show \rub, drop the student-only \stud
\ifdefined\showrub
	\renewcommand{\rub}[1]{\rubstyle{#1}}
	\renewcommand{\stud}[1]{}
\fi

% `pset -p` reads the points out of \prob{[20 pts (8, 12)]: title}, so the
% macro name is part of the contract; \probstyle takes the number and the
% title, and is where its look goes.
\newcounter{probcount}
\newcommand{\prob}[1]{\stepcounter{probcount}\probstyle{\arabic{probcount}}{#1}}
